Fix \(H\in\mathbb N\), instantiate \(T=H\), and fix a schedule satisfying Assumption \(\mathrm{ass:irreversible\text{-}endpoints}\).  We first establish the fixed-schedule one-wave equivalence.  For \(t\in\{1,\ldots,H\}\), put
\[
 a=p^{\mathrm{ext}}_{t-L},\qquad q=p_t.
\tag{1}
\]
Because \(t-L\leq t\), monotonicity and the endpoint extension give \(0\leq a\leq q\leq1\).

Suppose \(0<a\leq q<1\).  On \(\{U_i\leq a\}\), monotonicity gives \(U_i\leq p^{\mathrm{ext}}_{t-\ell}\) for every \(0\leq\ell\leq L\); hence all relevant exposures equal one.  On \(\{U_i>q\}\), one has \(U_i>p^{\mathrm{ext}}_{t-\ell}\) for every such \(\ell\); hence all relevant exposures equal zero.  With \(\theta_i=\sum_{\ell=0}^{L}b_{i,\ell}\), Assumption \(\mathrm{ass:additive\text{-}distributed\text{-}lag}\) therefore gives
\[
 Y^{\mathrm{obs}}_{i,t}=g_{i,t}+\theta_i
 \quad\text{on }\{U_i\leq a\},
 \qquad
 Y^{\mathrm{obs}}_{i,t}=g_{i,t}
 \quad\text{on }\{U_i>q\}.
\tag{2}
\]
The two events are disjoint because \(a\leq q\), and Assumption \(\mathrm{ass:nested\text{-}hash\text{-}randomization}\) assigns them probabilities \(a\) and \(1-q\).  Applying (2) to Definition \(\mathrm{def:interior\text{-}one\text{-}wave\text{-}ht\text{-}estimator}\) yields, unit by unit,
\[
 \mathbb E_U\{\phi_{i,t}(U_i)\}
 =(g_{i,t}+\theta_i)-g_{i,t}=\theta_i.
\tag{3}
\]
The statistic is measurable with respect to the one-wave data, all denominators are positive, and integrability follows because the potential outcomes form a fixed finite array under Assumption \(\mathrm{ass:fixed\text{-}potential\text{-}outcomes}\).  Definition \(\mathrm{def:outcome\text{-}wave\text{-}unbiased\text{-}rules}\) and (3) give
\[
 \widehat\theta_{\mathrm{one}}(t)\in
 \mathcal E_{\mathrm{wave}}(p,\{t\}),
 \qquad
 \mathbb E_U\!\left[\widehat\theta_{\mathrm{one}}(t)\right]
 =\frac1n\sum_{i\in I_n}\theta_i=\theta.
\tag{4}
\]

The same disjoint-event calculation gives
\[
 \mathbb E_U\{\phi_{i,t}(U_i)^2\}
 =\frac{(g_{i,t}+\theta_i)^2}{a}
  +\frac{g_{i,t}^2}{1-q}.
\tag{5}
\]
The scores \(U_i\) are independent across units, so the complete unit summands, but not repeated waves within a unit, are independent.  Subtracting the squared mean in (3), summing the unit variances in (5), and substituting (1) proves
\[
 \operatorname{Var}_U\!\left(\widehat\theta_{\mathrm{one}}(t)\right)
 =\frac1{n^2}\sum_{i\in I_n}
 \left\{
 \frac{(g_{i,t}+\theta_i)^2}{p^{\mathrm{ext}}_{t-L}}
 +\frac{g_{i,t}^2}{1-p_t}-\theta_i^2
 \right\}.
\tag{6}
\]

For the converse, suppose first that \(a=0\).  Compare the all-zero fixed population with the population in which \(b_{i,L}=1\) for every unit and all other \(g_{i,s}\) and \(b_{i,\ell}\) are zero.  At the sole observed time \(t\),
\[
 Z_{i,t-L}=\mathbf 1\{U_i\leq a\}=0
 \quad\text{almost surely}.
\tag{7}
\]
The two models have the same law for all data allowed by Definition \(\mathrm{def:outcome\text{-}wave\text{-}unbiased\text{-}rules}\), but their target values are zero and one.  Thus \(\mathcal E_{\mathrm{wave}}(p,\{t\})=\varnothing\).  If instead \(q=1\), compare the all-zero population with the population having \(b_{i,0}=1\), \(g_{i,t}=-1\), and all other relevant entries zero.  Since \(Z_{i,t}=1\) surely, the observed outcome at \(t\) is zero in both populations, while their targets are again zero and one.  Therefore \(\mathcal E_{\mathrm{wave}}(p,\{t\})=\varnothing\).  Since \(0\leq a\leq q\leq1\), these two aliases and (4) prove
\[
 \mathcal E_{\mathrm{wave}}(p,\{t\})\ne\varnothing
 \quad\Longleftrightarrow\quad
 0<p^{\mathrm{ext}}_{t-L}\leq p_t<1.
\tag{8}
\]

We now derive the bivariate frontier.  No zero-wave statistic can be uniformly unbiased: with no outcomes revealed, the all-zero population and the population with \(b_{i,0}=1\) have identical observed \((U,Z,p)\) but target values zero and one.  Thus every feasible wave count is at least one.

If \(L\geq1\) and \(H\leq L\), then every possible measured time satisfies \(t-L\leq0\).  Changing every \(b_{i,L}\) from zero to one leaves every revealed outcome unchanged almost surely, for every \(\mathcal S\), and changes \(\theta\) by one.  Hence no \(1\leq m\leq H\) is feasible.  If \(L=0\) and \(H=1\), the only possible wave is terminal and \(p_1=1\); the contemporaneous-treatment alias used after (7) proves infeasibility.

Let \(L\geq1\) and \(H=L+1\).  Proposition \(\mathrm{prop:elbow\text{-}wave\text{-}placement}\) proves the sharper placement certificate
\[
 \exists p:\ \mathcal E_{\mathrm{wave}}(p,\mathcal S)\ne\varnothing
 \quad\Longleftrightarrow\quad
 H\in\mathcal S\ \text{ and }\
 \mathcal S\cap\{1,\ldots,L\}\ne\varnothing.
\tag{9}
\]
In particular no one-wave set is feasible, whereas each pair \(\{s,H\}\), \(1\leq s\leq L\), is feasible for the strictly increasing schedule constructed there.  Thus the minimum wave count at this horizon is exactly two.  The pair \(\{1,H\}\) has the additional universal property stated in the target: for every fixed schedule with \(0<p_1\leq\cdots\leq p_{H-1}<1\), Definition \(\mathrm{def:panel\text{-}endpoint\text{-}overlap}\) holds, and Theorem \(\mathrm{thm:sharp\text{-}observed\text{-}panel\text{-}stage\text{-}frontier}\), applied at \(T=H=L+1\), shows that Definition \(\mathrm{def:panel\text{-}ht\text{-}estimator}\) is exactly unbiased.  Because that estimator uses only waves \(1,H\), it belongs to \(\mathcal E_{\mathrm{wave}}(p,\{1,H\})\).  This proves universal sufficiency of the canonical endpoint pair without claiming that it is the unique sufficient pair.

If \(L\geq1\) and \(H\geq L+2\), choose \(p_j=j/H\) and take \(t=L+1\).  Then
\[
 0<p_{t-L}=p_1\leq p_t=p_{L+1}<1
\]
because \(L+1\leq H-1\).  Equation (8) gives feasibility with one wave, and the zero-wave alias proves minimality.  If \(L=0\) and \(H\geq2\), the same schedule with \(t=1\) satisfies \(0<p_t<1\), so (8) again gives a one-wave estimator and minimality.

Finally, whenever \((H,m_0)\) is feasible and \(m_0\leq m\leq H\), enlarge the measured set to cardinality \(m\) and use the same statistic while ignoring the added outcomes.  This proves every upper portion of the two displayed feasible sets; the preceding aliases exclude every omitted pair.  Definition \(\mathrm{def:bivariate\text{-}rollout\text{-}wave\text{-}frontier}\) then gives the asserted formula for \(m_{\min}^{\mathrm{wave}}(H,L)\).  Under the convention that every rollout threshold is measured, Theorem \(\mathrm{thm:sharp\text{-}observed\text{-}panel\text{-}stage\text{-}frontier}\) supplies the diagonal identity \(T_{\min}^{\mathrm{panel}}(L)=\max\{2,L+1\}\).  Every proof above permits arbitrary fixed \(g\), so neither the bivariate frontier nor its placement refinement depends on \(r\).